% 本文件由 paper/prep_appendix_code.py 生成，勿手改

\begingroup\noindent\textbf{共用函数（四问依赖）}\par\endgroup

\subsection{读取附件 2 的负荷与光伏数据（\texttt{func_read_q2.m}，53 行）}
\VerbatimInput{code/func_read_q2.m}

\subsection{负荷与光伏的中心预测（\texttt{func_forecast_q2.m}，76 行）}
\VerbatimInput{code/func_forecast_q2.m}

\subsection{历史预测残差库（\texttt{func_resid_q2.m}，48 行）}
\VerbatimInput{code/func_resid_q2.m}

\subsection{按小时偏差校正（\texttt{func_bias_q2.m}，101 行）}
\VerbatimInput{code/func_bias_q2.m}

\subsection{写结果表 result*.xlsx（\texttt{func_write_q2.m}，159 行）}
\VerbatimInput{code/func_write_q2.m}

\begingroup\noindent\textbf{问题一}\par\endgroup

\subsection{入口：装配、求解、校验与落盘（\texttt{main_q1.m}，131 行）}
\VerbatimInput{code/main_q1.m}

\subsection{读取附件 1 的电价、负荷与光伏（\texttt{func_read_q1.m}，23 行）}
\VerbatimInput{code/func_read_q1.m}

\subsection{显式分流 MILP 的变量、目标与约束（\texttt{func_build_q1.m}，124 行）}
\VerbatimInput{code/func_build_q1.m}

\subsection{能量平衡与储能递推校验（\texttt{func_check_q1.m}，52 行）}
\VerbatimInput{code/func_check_q1.m}

\subsection{写 result1.xlsx 与表 1、表 2 数值（\texttt{func_write_q1.m}，59 行）}
\VerbatimInput{code/func_write_q1.m}

\begingroup\noindent\textbf{问题二}\par\endgroup

\subsection{入口：运行清单与逐日滚动（\texttt{main_q2c.m}，77 行）}
\VerbatimInput{code/main_q2c.m}

\subsection{7 日滚动主循环与断点续跑（\texttt{func_roll_q2c.m}，169 行）}
\VerbatimInput{code/func_roll_q2c.m}

\subsection{两阶段 SAA-MILP 的变量、目标与约束（\texttt{func_build_q2c.m}，193 行）}
\VerbatimInput{code/func_build_q2c.m}

\subsection{联合误差情景抽样（\texttt{func_scen_q2c.m}，50 行）}
\VerbatimInput{code/func_scen_q2c.m}

\subsection{按真实数据执行并结算（\texttt{func_exec_q2c.m}，74 行）}
\VerbatimInput{code/func_exec_q2c.m}

\begingroup\noindent\textbf{问题三}\par\endgroup

\subsection{入口：四阶段策略运行（\texttt{main_q3b.m}，115 行）}
\VerbatimInput{code/main_q3b.m}

\subsection{日内多阶段滚动主循环（\texttt{func_roll_q3b.m}，226 行）}
\VerbatimInput{code/func_roll_q3b.m}

\subsection{阶段调整与结算的 MILP 装配（\texttt{func_build_q3b.m}，267 行）}
\VerbatimInput{code/func_build_q3b.m}

\subsection{分阶段情景抽样（\texttt{func_scen_q3b.m}，73 行）}
\VerbatimInput{code/func_scen_q3b.m}

\subsection{阶段执行与拼接重放（\texttt{func_exec_q3b.m}，74 行）}
\VerbatimInput{code/func_exec_q3b.m}

\subsection{分阶段残差库（\texttt{func_resid_q3b.m}，75 行）}
\VerbatimInput{code/func_resid_q3b.m}

\subsection{附件 3 整点预报插值到 10 min（\texttt{func_interp_q3b.m}，23 行）}
\VerbatimInput{code/func_interp_q3b.m}

\subsection{读取附件 3 的分阶段预报（\texttt{func_read_q3b.m}，35 行）}
\VerbatimInput{code/func_read_q3b.m}

\begingroup\noindent\textbf{问题四}\par\endgroup

\subsection{入口：四种口径运行（\texttt{main_q4.m}，114 行）}
\VerbatimInput{code/main_q4.m}

\subsection{联合不确定性的滚动主循环（\texttt{func_roll_q4.m}，246 行）}
\VerbatimInput{code/func_roll_q4.m}

\subsection{电价中心预测与价格残差库（\texttt{func_price_q4.m}，105 行）}
\VerbatimInput{code/func_price_q4.m}

\subsection{价格与负荷光伏的联合情景（\texttt{func_scen_q4.m}，64 行）}
\VerbatimInput{code/func_scen_q4.m}

\subsection{电价预测入口（\texttt{func_forecast_q4.m}，21 行）}
\VerbatimInput{code/func_forecast_q4.m}
